	\section{动力学————牛顿运动定律}
	\subsection{牛顿运动定律}
	
	
	
	{\footnotesize 物理学史：牛顿（ I . Newton ）《自然哲学的数学原理》出版标志着经典力学体系的确立。}
	
	
	动力学是研究物体与物体之间的相互作用以及由于这种相互作用而引起的物体运动状态的变化规律，其基础是牛顿三大定律。
	\begin{theorem}[牛顿第一定律]
		任何物体都保持静止或沿一直线作匀速运动的状态，直到作用在它上面的力迫使它改变这种状态为止。
	\end{theorem}
	\begin{note}
	 	第一定律建立了惯性和力的确切概念。任何物体都具有惯性，因此第一定律又被叫做惯性定律( law of inertia )。惯性( inertia )：物体所具有的保持其原有运动状态不变的特性。
	
	牛顿第一定律动力学是研究物体之间的相互作用以及由于这种相互作用而引起的平动物体或质点运动状态的变化规律，也是其他定律的基础。
	
	惯性分类：质点运动的平动惯性，转动的转动惯性，热运动的热惯性，电磁运动的电磁惯性。
	
	牛顿第一定律对惯性参考系( inertial frame )都适用，即一个不受力作用的物体将保持其静止或匀速直线运动的状态不变。不遵守牛顿第一定律的参考系，称为非惯性参考系，简称非惯性系。
	\end{note}
	\begin{theorem}[牛顿第二定律]
		
		物体受到外力作用时，加速度的大小与外力的大小成正比，并与物体的质量成反比，加速度的方向与外力的方向相同。
	
\centering	$ \vec{F=ma } $
	
	单位：$kg ，m /s^2，N$ 
	
	动量( momentum )，用 $p$ 表示，即\\
	
	$ p=mv $
	
	第二定律实质上是\\
	
	$ \frac{dp}{dt}= F  $
	
	
	
	$ F=m\frac{d^{2}r}{dt^{2}} $
	
	称为质点的动力学方程。\\
	\end{theorem}

	\begin{note}
	牛顿第二定律注意：
	
	(1)瞬时性，第二定律定量地表述了物体的加速度与所受外力之间的瞬时关系，它们同时存在，同时改变，同时消失。
	
	(2)矢量性，力和加速度的叠加和作用效果遵从矢量三角形法则。
	
	(3)对应性，一个力对应一个加速度.
	\end{note}
	
	\begin{definition}[牛顿第三定律]两个物体之间的作用力和反作用力，在同一直线上，大小相等方向相反。


即
$ F_{ A B } = - F_{ B A }  $

若把
$F _ { AB } \text{和} F _ { BA }$中的一个叫做作用力（acting force），则另一个就叫反作用力（reacting force）。

第三定律揭示了力的特性：成对性，一致性，同时性。
\end{definition}
\subsection{力学中的常见力}

力学中常见力有万有引力、重力、弹性力、摩擦力等。

\begin{definition}[万有引力]
	
	
	牛顿万有引力定律（law of universal gravity）:质量分别为m1和m2的两个质点，相距为r时，引力为
	$ F = G \frac { m _ { 1 } m _ { 2 } } { r ^ { 2 } } $
	
	式中G叫做引力常量（gravitational constant），2018年国际推荐值为
	$ G = 6 . 6 7 4 3 0 \times 1 0 ^ { - 1 1 } N \cdot m ^ { 2 } / k g ^ { 2 } $
\end{definition}

式中的质量反映了物体的引力性质，叫做引力质量（gravitational mass），和反映物体惯性的惯性质量在意义上是不同的。

适用范围：两个质点，均匀质量球体。


\begin{definition}[重力]因地球吸引而使物体受到的力。\end{definition}


重力的方向，重力加速度的方向：竖直向下的．\\[0.01cm]

{\small	由于地球的自转效应，地球不是一个惯性系。从某个惯性系（如日心系）看来，地球上的物体将随地球绕地轴作圆周运动。物体所受地球的引力有一部分提供了向心力，剩余的分力才是重力重力的大小。
	
	重力探矿法:地球内某处存在大型矿藏，破坏了地球质量的对称分布，使该处的重力加速度值表现出异常，因此可通过重力加速度的测定来探矿。}\\[0.01cm]



\begin{definition}[弹力](elastic force):发生形变的物体，由于要恢复原状，对与它接触的物体会产生力的作用,简称弹力。
\end{definition}

弹力的主要三种表现形式：

(1)正压力：

正压力( normal pressure )或支持力( support force )：两个物体通过一定面积相互挤压。

{\small	如，屋架压在柱子上，重物放在桌面上。大小:挤压程度。方向:垂直于接触面而指向对方。}\\

(2)绳中的张力：

张力( tension )：把绳分为两部分，这种内部的两部分绳的相互拉力。

当绳没有加速度，或质量可以忽略时，绳上各点的张力都是相等的，而且就等于外力条件。\\

(3)弹簧的弹力：

弹力（恢复力）：当弹簧被拉伸或压缩时，它就会对与之相连的物体有弹力作用。在弹性限度内，其大小和形变成正比．以 F 表示弹力，以 x 表示形变亦即弹簧的长度变化，则

$	F =- kx  $

式中 k 叫做弹簧的劲度系数（ stiffness ），负号表示弹力的方向总是和弹簧位移的方向相反，弹力总是指向要恢复它原长的方向。\\

4.摩擦力

摩擦力( friction force ):两个相互接触的物体在沿接触面相对运动时，或者有相对运动的趋势时，在接触面之间产生一对阻碍相对运动的力。\\

(1)静摩擦力( static friction force):相互接触的两个物体在外力作用下，虽有相对运动的趋势，但并不产生相对运动。\\

每个物体所受静摩擦力的方向与该物体相对于另一物体的运动趋势方向相反。静摩擦力的大小视外力的大小而定，介乎0和最大静摩擦力 F 之间。最大静摩擦力正比于正压力 F ，即\\

$	F_{s}=\mu_{s} F_{N}   $

静摩擦因数$\mu_{s}$( static friction factor )，与接触面的材料和表面情况有关。\\

(2)动摩擦力( dynamic friction force ):当外力超过最大静摩擦力时，物体间将会产生相对运动的摩擦力。\\

动摩擦力 F 也与正压力 F 成正比，\\

$	F_{k} = \mu_{k} F_{N}  $
动摩擦因数$\mu_{k}$( dynamic friction factor )和相互接触的两物体的材料，表面情况，物体的相对速度有关。在大多数情况下，它随速度的增大而减小。\\

对于给定的一对接触面来说， $\mu_{s}> \mu_{k}$。\\

(3)湿摩擦力:固体在流体中运动时，沿着接触面出现的摩擦力。\\

1：在固体与流体相对运动速率不大时\\
$	F_{d} = -kv $

2：在固体与流体相对运动速率较大时\\
$	F_{d} =-kv\overrightarrow{v} $

湿摩擦力远小于干摩擦力！
\subsection{基本相互作用}	
自然界物体之间的相互作用力，可归结为四种相互作用：\\

引力相互作用:物体之间相互吸引的作用，这种力存在于一切物体之间，其强度随距离的增大而减弱。


电磁相互作用:电荷之间的相互作用，包括同种电荷之间的排斥和异种电荷之间的吸引，以及磁体之间的相互作用。


强相互作用:原子核内部将质子粘结在一起的力，仅在原子核内部起作用，外部没有任何表现。


弱相互作用:基本粒子之间的一种特殊作用，其强度比强相互作用和电磁作用都弱，作用范围也比强相互作用更小。



\section{非惯性系和惯性力}

\begin{definition}[非惯性系]
    惯性系：相对一个惯性系作匀速直线运动的一切参考系。\\
    非惯性系:相对地面参考系（惯性系）作加速运动的参考系。\\
\end{definition}


\begin{example}[非惯性系理解]
    以加速启动的列车为参考系，在列车车厢内的人看到小车的情况就大不一样，小车是向
    列车车尾方向作加速运动的．小车受力的情况没有变化，合力仍然是0，却有了加速度，
    这是违反牛顿运动定律的．因此，加速运动的列车是个非惯性系，相对于它，牛顿运动
    定律不再成立．
\end{example}


\begin{definition}[惯性力]
    设有一质点，质量为$m$ ，所受的合外力为 $\vec{F}$ ，它相对于惯性系的加速度为 a ，根
据牛顿第二定律，有\\
\centering$\vec{F} = m\vec{a}$\\
设想另一参考系，相对于惯性系以加速度 $\vec{a}$ 。作直线运动，在此参考系中，质点的加速度为 $\vec{a'}$，
由运动的相对性,可知\\
 $\vec{a = a '+a_0}$\\
将此式代人牛顿运动定律可得\\
 $\vec{F = m ( a '+ a)= ma '+ ma}$\\
将上式改写成\\
 $\vec{F +(- ma_0 )= ma '}$\\
此式说明，在非惯性系观察物体运动时，除了实际的外力之外，物体还受到一个大小和方向由$﹣ m\vec{a}$ 。
表示的力．这个力称为惯性力．这样就可以形式上应用牛顿第二定律去分析问题．
\par
在直线加速的非惯性系中，质点$m$受到的惯性力$\vec{F_惯}$：\\
$\vec{F_惯}=-m\vec{a_0}$\\


\begin{note}
    \textcolor{red}{\kaishu 公式中负号为与加速度$\vec{a_0}$方向相反.$\vec{a_0}$为系统加速度.}
\end{note}


在匀速转动的参考系中，质点$m$受到的惯性力$\vec{F_惯}$：\\
$\vec{F_惯}=-mw^2r\vec{e_r}$\\
因此：若质点静止在匀速转动的参考系中：\\
$\sum \vec{F_外}+\vec{F_惯}=\vec{0}$ \\


\begin{note}
    \textcolor{red}{\kaishu 此力称为惯性离心力.}
\end{note}


\par
\par


\begin{supplement}[科里奥利加速度]
    \textcolor{blue}{\kaishu 物体在做牵连运动的同时,沿旋转半径做相对运动,由牵连运动和相对运动交互
    耦合而形成的加速度称为科里奥利加速度.\\}
    \centering$\vec{F_科}=2m\vec{v}\times\vec{\omega } $\\
    \begin{note}
        \textcolor{red}{\kaishu 1、该公式中角速度$\vec{\omega }$为矢量，方向可由右手螺旋定则确定.\\
        2、公式中$\vec{v}$为物体相对匀速转动非惯性系的相对速度.}
    \end{note}

\end{supplement}

\end{definition}


